\documentclass[preprint,12pt]{elsarticle}

\usepackage{amsmath,amssymb}
\usepackage{graphicx}
\usepackage{booktabs}
\usepackage{array}
\usepackage{tabularx}
\usepackage[hidelinks]{hyperref}

\newcolumntype{Y}{>{\raggedright\arraybackslash}X}
\newcolumntype{L}[1]{>{\raggedright\arraybackslash}p{#1}}

% A [GAP] placeholder for a table or figure whose content is not yet available.
\newcommand{\gapbox}[1]{\par\medskip\noindent\fbox{\parbox{\dimexpr\linewidth-2\fboxsep-2\fboxrule\relax}{#1}}\par\medskip}

% Vancouver reference style: requires vancouver.bst (CTAN package "vancouver") on the TeX path.
% References are numbered in order of first citation; sort&compress turns [3,1,2] into [1-3].
\biboptions{sort&compress}

\journal{Artificial Intelligence in Medicine}

\begin{document}

\begin{frontmatter}

\title{An interpretable block-term tensor network for incident atrial fibrillation prediction from the sinus-rhythm 12-lead electrocardiogram and clinical risk factors\tnoteref{rt}}
\tnotetext[rt]{Running title: Interpretable block-term tensor network for incident AF}

\author{Authors, affiliations, ORCID identifiers and corresponding author: \textbf{[GAP-AUTH]}. Filled by the author list, author order, affiliations, ORCID identifiers and corresponding-author contact that the authors supply.}

\begin{abstract}
\noindent \textbf{Background:} Deep learning models predict incident atrial fibrillation (AF) from a sinus-rhythm 12-lead electrocardiogram (ECG), but their evidence is not designed to be read and their explanations are seldom tested for faithfulness.

\noindent \textbf{Methods:} This retrospective study used each patient's earliest sinus-rhythm ECG in MIMIC-IV-ECG, with the first hospital-coded AF or flutter after a 30-day blanking window as the outcome. A block-term tensor network (BTTN) projects a log-spectrogram tensor onto 12 learned frequency-by-lead blocks before a nonlinear stage trained under a Cox loss, and BTTN + clinical adds age, sex and seven ECG measurements as a linear term outside the block layer. Models were compared by the area under the receiver operating characteristic curve (AUROC), with patient-bootstrap 95\% confidence intervals (CIs), in repeated five-fold cross-validation. Five faithfulness analyses of the block signatures are described.

\noindent \textbf{Results:} Of 151,941 eligible patients, 1,855 were one-year cases and 41,770 were controls. BTTN + clinical had the highest point AUROC of seven models, 0.7776 (95\% CI 0.7664 to 0.7888), against 0.7597 (0.7488 to 0.7713) for a clinical model and 0.735 (0.7226 to 0.7469) for a convolutional neural network (CNN). ECG-only BTTN reached 0.7397 (0.7275 to 0.7517). The BTTN + clinical survival model had a higher point AUROC than a clinical Cox model at all four horizons, reaching 0.745 against 0.731 at five years. Observed one-year AF was 0.55\% in the lowest decile of BTTN + clinical risk and 11.4\% in the highest.

\noindent \textbf{Conclusions:} BTTN + clinical reached the highest point AUROC while all ECG evidence passed through readable frequency-by-lead block projections, and one survival model covered horizons from 90 days to five years. Faithfulness results are held at \textbf{[GAP-F1 to GAP-F5]}.

\end{abstract}

\begin{highlights}
\item A block-term tensor network reads the ECG through 12 learned frequency-by-lead blocks
\item Age, sex and seven ECG measurements enter as a linear term outside the block layer
\item With clinical fusion, one-year AF point AUROC was 0.7776, highest of seven models
\item ECG-only network reached AUROC 0.7397 against 0.735 for a CNN on the same tensor
\item The fused survival model exceeded the clinical Cox point AUROC at all four horizons
\item Observed one-year AF was 0.55\% in the lowest risk decile and 11.4\% in the highest
\item Five faithfulness analyses test whether the block signatures track the model
\end{highlights}

\begin{keyword}
tensor decomposition \sep deep learning \sep survival analysis \sep faithfulness evaluation \sep risk stratification \sep MIMIC-IV
\end{keyword}

\end{frontmatter}

\section{Introduction}

\subsection{AF burden and the case for risk-targeted screening}

The incidence and prevalence of atrial fibrillation (AF) rise with population ageing, and reviews of AF epidemiology treat risk prediction, together with modifiable risk factors and prognosis, as an active field that includes eHealth and artificial intelligence (AI) developments~\cite{kornej2020}. Targeted screening depends on risk prediction. Noseworthy and colleagues tested this in a prospective non-randomised trial of 1,003 patients who were stratified by an electrocardiogram (ECG)-based risk score and monitored for up to 30 days~\cite{noseworthy2022}. AF was detected in 7.6\% of high-risk and 1.6\% of low-risk patients, and in the high-risk group AI-guided screening raised detection from 3.6\% under usual care to 10.6\%~\cite{noseworthy2022}. A score that ranks patients well therefore raises the yield of each monitored patient.

The 12-lead ECG is a natural input for such a score. It is recorded routinely, and sinus-rhythm markers such as P-wave duration are associated with later AF~\cite{nielsen2015}. This paper asks how to build a risk model from one such recording whose ECG evidence can be read and checked.

\subsection{Clinical risk scores and AI-ECG models for incident AF}

Clinical risk models set a high standard. The Framingham AF risk score reached a C statistic of 0.78~\cite{schnabel2009}, and the CHARGE-AF model, developed in pooled cohorts of 18,556 participants, reached a five-year C statistic of 0.765 that adding ECG variables did not meaningfully change~\cite{alonso2013}. Deep models applied to the sinus-rhythm ECG have reached similar or higher discrimination for related tasks. A convolutional neural network (CNN) identified patients with AF during sinus rhythm with an area under the receiver operating characteristic curve (AUROC) of 0.87, a task of identifying concurrent or recent AF~\cite{attia2019}. A deep network predicted new-onset AF within one year with an AUROC of 0.85~\cite{raghunath2021}. A model trained on 669,782 sinus-rhythm ECGs predicted five-year incident AF or flutter with an AUROC of 0.78 internally and 0.77 in MIMIC-IV, used only as an external test set~\cite{jabbour2024}, and a model trained on ECGs of United States veterans flagged AF within 31 days with an AUROC of 0.86~\cite{yuan2023}.

ECG and clinical information carry partly separate signal. Khurshid and colleagues reported five-year AUROCs of 0.802, 0.823 and 0.838 for CHARGE-AF, an ECG model and their combination at one hospital, and 0.732, 0.705 and 0.746 in UK Biobank, where limited follow-up led them to recalibrate the models to two-year risk, and called the two approaches complementary~\cite{khurshid2022}. Christopoulos and colleagues found C statistics of 0.69, 0.69 and 0.72 for AI-ECG, CHARGE-AF and their combination~\cite{christopoulos2020}. Biton and colleagues reported that combining engineered ECG features, learned representations and clinical data raised five-year AF AUROC to 0.909, against 0.839 for the best single modality~\cite{biton2021}. A clinical risk-factor model is therefore a strong comparator, and fusion of ECG and clinical information is the design these studies favour.

\subsection{The interpretability gap}

The incident-AF models above are deep networks trained end to end on the ECG~\cite{attia2019,raghunath2021,jabbour2024}, and their internal evidence is not designed to be read. Where the sources reviewed offer an explanation for an ECG model, it is computed after training: local interpretable model-agnostic explanations (LIME) for a 12-lead ECG CNN~\cite{hughes2021}, gradient-weighted class activation mapping (Grad-CAM) heat maps for an ECG network~\cite{storas2023}, and Shapley additive explanation (SHAP) values for AF risk models built on clinical and ECG-derived variables~\cite{matsuo2026}. For the SHAP explanations, the abstract reports no test of faithfulness to the model~\cite{matsuo2026}, and physicians who reviewed the Grad-CAM heat maps found that they gave no meaningful information~\cite{storas2023}. In the literature searched for this study, no ECG study with a quantitative faithfulness evaluation was found.

Post-hoc attributions have known failure modes. Some saliency methods produce maps that are independent of both the model parameters and the training data~\cite{adebayo2018}, and under removal and retraining many popular importance estimates did no better than random importance~\cite{hooker2019}. Rudin argued that high-stakes decisions such as those in health care should use models that are interpretable by design~\cite{rudin2019}. Ghassemi and colleagues took a more sceptical view: current explainability methods are unlikely to deliver the trust and transparency expected for decisions about individual patients, and rigorous internal and external validation should take their place~\cite{ghassemi2021}.

Both positions imply that a model called interpretable needs a test of how far its readout follows its computation. The model here is interpretable by construction in a stated and limited sense, and its signatures are submitted to faithfulness tests.

\subsection{Aim and contributions}

The research question is whether an interpretable block-term tensor network (BTTN), whose ECG evidence passes through readable frequency-and-lead block signatures and is fused with age, sex and ECG intervals, can predict incident AF from one sinus-rhythm 12-lead ECG in MIMIC-IV-ECG on hospital-coded diagnosis labels more accurately than a clinical risk-factor model and a CNN, and whether those signatures are faithful. The paper makes four contributions. First, it defines BTTN, a block-term factorisation of the ECG spectral tensor into frequency and lead factors, trained end to end with a nonlinear temporal stage and head under a Cox partial-likelihood loss. Second, it adds clinical risk factors through a linear clinical term on the log-risk that bypasses the block layer. Third, it compares BTTN with a clinical model, a CNN and a linear block-term regressor on one MIMIC-IV-ECG cohort with International Classification of Diseases (ICD) outcome codes and shared folds. Fourth, it specifies a five-part faithfulness protocol. The name tensor network refers to the first layer, a set of block terms of the input tensor, and the model uses no chained tensor contractions.

\section{Related work}

\subsection{Block-term tensor models and tensor regression}

A block-term decomposition writes a higher-order tensor as a sum of terms with specified mode-n ranks, and it contains the canonical polyadic (CP) and Tucker decompositions as special cases~\cite{delathauwer2008}. Supervised multilinear regression builds on this structure. Higher-order partial least squares predicts a tensor from a tensor through a sum of Tucker blocks fitted by deflation~\cite{zhao2013}, and a fast variant reached comparable accuracy at lower cost~\cite{camarrone2019}. Block-term tensor regression (BTTR) with automatic component extraction predicted single-finger trajectories from intracranial recordings~\cite{faes2024}, and an extended form based on recursive Tucker decomposition predicted finger movement and coactivation~\cite{faes2022}. These models are linear in the input, and BTTR serves in this study as the linear comparator.

Block terms have also been read as physiological components. Hunyadi and colleagues decomposed wavelet-transformed EEG with rank-(Lr, Lr, 1) block terms to capture seizure patterns~\cite{hunyadi2014}, and Goulart and colleagues estimated block-term structure jointly with its parameters to extract atrial activity from the ECG during AF~\cite{goulart2020}. Both are unsupervised decompositions. Kim and colleagues fitted a supervised CP factorisation of diagnosis-by-prescription tensors on MIMIC-III that gave readable phenotypes and predicted intensive care mortality~\cite{kim2017}. Van Den Kerchove and colleagues proposed block-term tensor discriminant analysis, which extracts supervised block-term features from EEG for a separate classifier~\cite{vandenkerchove2025}.

In the sources reviewed, block terms inside neural networks served to compress weights. Block-term decompositions reduced the parameters of recurrent networks~\cite{ye2018} and of fully connected layers~\cite{li2017btnets}, and Ye and colleagues approximated the weight matrices of CNNs and recurrent networks with low-rank block-term tensors under the name block-term tensor neural networks~\cite{ye2020}. BTTN shares the name and differs in mechanism. It factorises the input tensor into frequency and lead factors that lie on the predictive path, and its block terms define how the input is read. Dooms and colleagues combined multilinear tensor-network structure with deep network capacity for weight-based interpretability of image models~\cite{dooms2025}, without a clinical outcome or physiological signatures. To our knowledge, and within the literature searched for this study, no prior model trains a block-term input factorisation end to end with a nonlinear head and a survival objective for a clinical outcome.

\subsection{Interpretable by design and faithfulness evaluation}

Interpretable-by-design models place a readable representation on the predictive path. Concept bottleneck models first predict human-specified concepts and then the label, and they reached accuracy competitive with end-to-end models on x-ray grading and bird identification while allowing concepts to be corrected at test time~\cite{koh2020}. B-cos networks make the whole network reduce to one input-dependent linear transform that summarises its computation, with ImageNet accuracy similar to standard architectures~\cite{bohle2022}. On the tasks studied, a readable intermediate layer was compatible with competitive accuracy.

A readable layer that feeds a flexible predictor has a known weakness. Mahinpei and colleagues showed that learned concept representations leak information beyond the predefined concepts and that common mitigations only partly remove it~\cite{mahinpei2021}. Havasi and colleagues showed that training the label predictor on soft concept probabilities, which improves accuracy, lets information leak, which corrupts concept predictions and weakens intervention~\cite{havasi2022}. For BTTN the corresponding question is whether the nonlinear temporal stage and head use information that the readable block projections do not expose. A test of this question has to read the block projections without the head, so that the head's capacity cannot enter the readout.

Faithfulness is the property such tests measure. Jacovi and Goldberg separated faithfulness, the degree to which an interpretation reflects the model's reasoning, from plausibility to a human reader, and argued for a graded notion of faithfulness~\cite{jacovi2020}. Two families of tests follow from this literature. Randomisation tests check whether an explanation changes when model parameters or training labels are randomised~\cite{adebayo2018}. Removal tests check whether deleting the features that an explanation ranks as important degrades a model more than deleting others~\cite{hooker2019}. Integrated Gradients, an attribution method built on axioms of sensitivity and implementation invariance~\cite{sundararajan2017}, represents the explanations computed after training that these tests were designed to check. In the sources reviewed, no study combines a supervised, survival-trained block-term input factorisation with measured faithfulness for a clinical outcome. Section 3 defines the cohort, the model and the tests that address this combination.

\section{Methods}

\subsection{Study design, data sources and reporting}

This is a retrospective development and internal validation study of a prognostic model with repeated patient-level cross-validation. ECGs came from MIMIC-IV-ECG version 1.0, which holds about 800,000 diagnostic 12-lead ECGs from nearly 160,000 patients, each 10 seconds long and sampled at 500 Hz, matched to the MIMIC-IV clinical database~\cite{gow2023}. Diagnoses, admissions and patient records came from the hospital module of MIMIC-IV, drawn from records of Beth Israel Deaconess Medical Center~\cite{johnson2023}. Both resources are hosted on PhysioNet~\cite{goldberger2000}. Each patient contributed one ECG, the index ECG defined below, so every MIMIC-IV analysis is at the patient level and a resample of rows is a resample of patients. Reporting follows the TRIPOD+AI statement for prediction models that use regression or machine learning~\cite{collins2024}. The checklist is \textbf{[GAP-TRIPOD]}, filled by the TRIPOD+AI checklist that the authors complete against the final manuscript. The released analysis code fixes every setting stated in this section (Data availability).

\subsection{Cohort and outcome definition}

Each patient's index ECG was the earliest ECG whose machine report contained the word sinus (including sinus tachycardia, bradycardia and arrhythmia) and mentioned neither AF or atrial flutter nor a paced, junctional, atrial tachycardia, supraventricular tachycardia, wandering or ectopic atrial rhythm. The outcome was the first admission carrying ICD-10 code I48 (AF and flutter) or ICD-9 code 427.31 or 427.32, dated by admission. Patients with an outcome or an AF-reading ECG no later than 30 days after the index ECG were excluded as prevalent or concurrent AF.

In the one-year binary cohort, cases had an outcome more than 30 and at most 365 days after the index ECG. Controls had no outcome code at any time and at least 365 days of follow-up, ending at the later of the last ECG and death. An AF-reading ECG after the blanking window did not exclude a control, since the outcome was defined by hospital codes. In the survival cohort, an event was any outcome after the blanking window, and patients without one were censored at the end of follow-up of at least 365 days. Each cohort file holds at most 4,000 cases and 16,000 controls, sampled at random, and the one-year comparison kept all cases with four controls per case.

Nine covariates entered the clinical models and the clinical term: age (the MIMIC-IV anchor age, not adjusted to the date of the index ECG), sex, and seven machine measurements of the index ECG, namely RR interval, P-wave duration, an established predictor of AF~\cite{nielsen2015}, QRS duration, QT interval, and the P, QRS and T axes. The three durations were computed from machine-measured onset and end points. The machine-measurement missing codes 29999 and 32767 were not recoded before these differences were formed or the axes entered the models, so imputation applies only to values absent from the file. In the one-year cohort file, 819 of 17,855 rows (4.6\%, derived) carry a P-wave duration between 29,547 and 29,838 ms, and 37 axis values are 29999 or 32767 (counts of this file).

\subsection{ECG preprocessing and spectral tensor}

Each lead was filtered with a second-order Butterworth high-pass filter at 0.5 Hz and a notch filter at 50 Hz (quality factor 30), both applied forward and backward, and denoised by soft thresholding of a level-3 sym4 wavelet decomposition. The signal was decimated from 500 to 250 Hz with zero-phase anti-aliasing, centred, and scaled per lead to unit maximum absolute amplitude. For each lead, a magnitude spectrogram was computed with a 128-sample window and a 32-sample hop, frequencies up to 40 Hz were kept, magnitudes were transformed as log(1 + x), and the time axis was resampled to 40 frames. Each recording thus gives a frequency x time x lead tensor, standardised by its own mean and standard deviation, with bins about 1.95 Hz apart (derived from these settings). A spectrogram tolerates shifts in beat timing, so the tensor needs no beat detection or alignment. Time-frequency 12-lead representations have served as CNN inputs for arrhythmia detection~\cite{jeong2021}.

\subsection{BTTN model and training}

BTTN has four parts: a block-term layer, a temporal stage, a head and a survival loss (Figure 1). Let $\mathcal{X} \in \mathbb{R}^{F \times T \times L}$ be the spectral tensor with $F$ frequency bins, $T = 40$ time frames and $L = 12$ leads. The block-term layer has $K = 12$ blocks. Block $k$ holds a frequency factor $A_k \in \mathbb{R}^{F \times 3}$ and a lead factor $C_k \in \mathbb{R}^{L \times 2}$ and projects the tensor onto them with mode products in the notation of Kolda and Bader~\cite{kolda2009}:

\[\mathcal{Z}_k = \mathcal{X} \times_1 A_k^\top \times_3 C_k^\top \in \mathbb{R}^{3 \times T \times 2}.\]

Each block yields six channels over time, and the 12 blocks yield 72 channels. The layer has no time factor and no learned core tensor. It is a set of block terms in the sense of De Lathauwer~\cite{delathauwer2008} restricted to the frequency and lead modes, and time is left to the temporal stage. Every later computation receives only these 72 channels, so the frequency and lead directions available to the model are the columns of $A_k$ and $C_k$.

The temporal stage applies grouped one-dimensional convolutions along time to each channel, with kernel sizes 3, 7 and 15 and six filters per channel and kernel size, followed by batch normalisation and a rectified linear unit. The sequences are pooled over time by a learned softmax attention and by max pooling, concatenated. The head applies batch normalisation and dropout (rate 0.4), then a multilayer perceptron with hidden layers of 256 and 128 units, and outputs a scalar ECG log-risk.

BTTN was trained with the negative Cox partial log-likelihood~\cite{cox1972}, as in DeepSurv~\cite{katzman2018}, computed within each mini-batch so that a patient's risk set is the batch members with an equal or longer time, as in sampled risk sets~\cite{kvamme2019}. Cases contributed their time to the outcome and controls their follow-up time. Inputs were standardised per lead with training-fold statistics. Optimisation used Adam (learning rate 2 x 10$^{-3}$, weight decay 10$^{-4}$), a cosine schedule, batches of 512 and at most 120 epochs, with early stopping on the validation AUROC of a 15\% inner split (patience 16). Three models with different seeds were trained per fold and their risks averaged.

Two readouts are interpretable without the head. The block scores are the 72 projected channels averaged over time, a linear function of the input tensor that is multilinear in the factors. The block signature of block $k$ is the frequency x lead map $A_k C_k^\top$, formed from the first two columns of $A_k$ and both columns of $C_k$. It summarises two of the six channels of each block, and the complete readout of block $k$ is the set of six rank-one maps $a_{k,p} c_{k,s}^\top$. The temporal stage and the head are nonlinear and are not interpretable by design, which is why Section 3.8 tests faithfulness.

For BTTN + SSL and BTTN + SSL + clinical, the block-term layer and temporal stage were first pretrained by self-supervised learning (SSL) with a SimCLR-style contrastive objective~\cite{chen2020simclr}, as adapted to cardiac signals~\cite{kiyasseh2021} and 12-lead ECGs~\cite{mehari2022}. Pretraining used the NT-Xent loss (temperature 0.2) on 30,000 unlabelled ECGs sampled one per patient from MIMIC-IV-ECG for 60 epochs, with views formed by time and frequency masking, lead dropout, additive noise and a small roll along time (settings in Table S1).
% Figure 1 is not rendered: paper/figures/fig1_schematic.pdf does not exist. When it is rendered,
% replace this box with a figure environment that includes ../figures/fig1_schematic.pdf and carries this caption text.
\refstepcounter{figure}
\gapbox{\textbf{[GAP-FIG1]} \textbf{Figure 1.} Model schematic: the spectral tensor, the block factors $A_k$ and $C_k$ without a core tensor, the grouped temporal convolutions with attention and max pooling, the multilayer perceptron head, the additive clinical term on the nine covariates and the Cox loss. Filled by rendering the corrected schematic plotting script delivered with the code.}
\label{fig:bttn-schematic}

\subsection{Clinical fusion term}

BTTN + clinical adds a linear clinical term to the ECG log-risk,

\[\eta = h(\mathcal{X}) + w^\top c,\]

where $h(\mathcal{X})$ is the head output and $c$ is the vector of the nine covariates of Section 3.2, standardised with training-fold means and standard deviations. Absent age was set to the training mean and absent sex was coded as not male (2 of 17,855 rows each in the one-year cohort file), and values carrying a missing code enter as numbers. The term has no bias. Its weights $w$ were initialised from a class-balanced logistic regression of the one-year label on the covariates of the training patients and excluded from weight decay, so that weight decay does not shrink them. Because the clinical term is added after the head, it enters neither the block-term layer nor the block scores, and the block scores of BTTN + clinical are computed exactly as for BTTN. The factor values are learned jointly with the clinical term, so they can differ between the two models while the structure stays the same.

\subsection{Comparators}

Six comparators shared the cohort, folds and spectral tensor of BTTN + clinical. The clinical model was an L2-penalised logistic regression on the nine covariates with median imputation of absent values, standardisation and balanced class weights, in the role clinical risk scores play as AF baselines~\cite{schnabel2009,alonso2013,khurshid2022}. Linear BTTR was block-term tensor regression with automatic component extraction~\cite{faes2024}, using the bttrlib implementation released with the analysis code, fitted to the label coded as -1 and +1 with up to eight blocks (search grid in Table S1).

The CNN was a compact network of three convolutional layers, used as a black-box comparator on the same input, in the model class of AI-ECG work~\cite{attia2019,raghunath2021}. It received the tensor with frequency and lead flattened into channels, standardised per channel. Its layers had 32, 64 and 64 filters with kernel sizes 7, 5 and 3, batch normalisation, rectified linear units, dropout of 0.3, 0.4 and 0.5, max pooling after the first two and average pooling after the third. Training used binary cross-entropy weighted by the control-to-case ratio, Adam and early stopping on validation AUROC on one inner 85/15 split shared by a three-model ensemble per fold (settings in Table S1). ECG-only BTTN, BTTN + SSL and BTTN + SSL + clinical complete the set. In the multi-horizon analysis the clinical comparator was a Cox model with an L2 penalty of 0.1 on the same covariates and median imputation, the clinical Cox model.

\subsection{Evaluation protocol and statistics}

In the one-year comparison, stratified five-fold cross-validation was repeated with five seeds, each drawing new folds, and each patient's out-of-fold scores were averaged over seeds. AUROC, the probability that a random case is ranked above a random control~\cite{hanley1982}, was computed in 2,000 patient-bootstrap resamples of the averaged scores and is reported as the mean over resamples with the 2.5th and 97.5th percentiles as the 95\% CI~\cite{carpenter2000}. Comparisons are reported as point estimates with CIs, and the paired patient-bootstrap AUROC differences (2,000 resamples) against the clinical model and the CNN are held at \textbf{[GAP-C4]} in Section 4.2.

The multi-horizon analysis used the survival cohort in a single five-fold run stratified by event. One BTTN and one BTTN + clinical were trained on continuous time, each fold a three-model ensemble with early stopping on validation AUROC for an event within one year. AUROC was read from the same out-of-fold risk at four horizons, with cases having an event by the horizon, controls an event or censoring after it, and earlier censoring excluded. The 90-day horizon counts outcomes from day 31 to day 120. CIs for the BTTN models came from 2,000 patient-bootstrap resamples.

Calibration and clinical utility were assessed at one year on the multi-horizon scores, with probabilities from two-fold cross-fitted Platt scaling: reliability curve and Brier score~\cite{vancalster2019}, decision curve~\cite{vickers2006}, and sensitivity and specificity when flagging the top 5\%, 10\%, 20\% and 30\% of patients. Positive predictive value (PPV) and the number needed to monitor (1 / PPV) were recomputed at the one-year incidence of the eligible cohort, 1,855 / (1,855 + 41,770). Observed one-year AF rates were tabulated by decile of BTTN + clinical risk.

\subsection{Faithfulness analyses}

Faithfulness is the degree to which a readout of BTTN agrees with the model's own computation, graded rather than binary~\cite{jacovi2020}. All five analyses use the one-year comparison set and ECG-only BTTN as in Section 3.4, as single models in the probes, intervention test and block-structure ablation, with five folds and three seeds where cross-validated.

\textit{Control probes.} A logistic regression fitted on training-patient block scores was evaluated on the test fold. Block scores exclude the temporal stage and head, so the probe cannot use the head's capacity, the leakage route of concept models~\cite{mahinpei2021,havasi2022}. As in randomisation checks~\cite{adebayo2018}, the same probe was fitted to four controls: the raw flattened tensor, the 64-dimensional penultimate CNN features, a fixed random projection to the block-score dimension, and block scores of the label-blind self-supervised encoder. Probes are reported as AUROC with CI and as recovery, probe AUROC divided by full BTTN AUROC on the same folds.

\textit{Intervention test.} On a held-out fold, frequency bins and leads were ranked by the AUROC drop when each was set to zero alone. The top-n and, as a control, bottom-n axes were then zeroed cumulatively and full-model AUROC recorded on the same fold without retraining, unlike the remove-and-retrain benchmark~\cite{hooker2019}.

\textit{Block-structure ablation.} The block-term layer was replaced by 12 rank-one CP blocks (12 channels) or by a full-rank frequency x lead projection with the same 72 channels, downstream unchanged, recording AUROC and probe recovery.

\textit{Block-count sweep.} $K$ was set to 4, 8, 12, 16 and 24, reporting AUROC mean and standard deviation over seeds.

\textit{Substrate map.} On a held-out 20\% split, a three-model ensemble gave occlusion importance per frequency bin and per lead and the mean absolute block signature, in its two-channel form (Section 3.4).

\subsection{External analyses}

Two independent datasets without dated outcomes test transfer of the method through AF identification or detection, tasks that differ from incident prediction. CODE-15, a stratified 15\% sample of the CODE study of the Telehealth Network of Minas Gerais, holds 345,779 exams from 233,770 patients with an AF flag per exam~\cite{ribeiro2021code15,ribeiro2020}. With one exam per patient and at most 6,000 per class, controls were patients never flagged with AF, and cases were non-AF exams of patients with an AF exam elsewhere (identification) or AF exams (detection). PTB-XL holds 21,837 records from 18,885 patients~\cite{wagner2020}, from which a class-balanced task detecting AF or flutter against sinus rhythm was sampled by record, giving record-level folds. Undecimated tracings at 400 Hz (CODE-15) and 100 Hz (PTB-XL) give bins about 3.1 Hz and 0.78 Hz apart (derived), and at 100 Hz the 50 Hz notch lies at the Nyquist frequency. A single BTTN with binary cross-entropy loss was evaluated by five-fold cross-validated AUROC, and occlusion importances per frequency bin and lead were computed in-sample on a model fitted to all records.

\section{Results}

\subsection{Cohort}

Of 151,941 patients with an eligible sinus-rhythm index ECG, 14,469 met the criterion for prevalent or concurrent AF, 88,945 had no outcome and under 365 days of follow-up, and 7,481 had their first outcome after one year (Table 1). These counts overlap and are not subtracted in sequence. The eligible one-year cohort comprised 1,855 cases and 41,770 controls (43,625 patients), and the survival cohort 8,775 patients with an incident outcome and the same controls (50,545 patients). Incident outcomes numbered 1,855 by one year, 2,828 by two years and 5,262 by five years. Controls had a median follow-up of 1,601 days (interquartile range 880 to 2,559). Of the one-year controls, 98.1\% had at least one later ECG and 96.5\% had a latest later ECG that did not read AF.

The one-year cohort file held all 1,855 cases and 16,000 controls. The one-year comparison set comprised the 1,855 cases and 7,420 controls, 9,275 patients at a prevalence of 0.20 (derived from four controls per case). The survival cohort file held 4,000 events and 16,000 censored patients. Baseline characteristics by outcome: \textbf{[GAP-T1]} (Table 1), filled by descriptive statistics of age, sex and the seven ECG measurements by outcome in the one-year and survival cohort files.

\begin{table}[htbp]
\centering
\small
\caption{Cohort derivation and follow-up for ICD-coded outcomes.}
\label{tab:cohort}
\begin{tabularx}{\textwidth}{@{}L{0.42\textwidth}YY@{}}
\toprule
Quantity & One-year binary cohort & Survival cohort \\
\midrule
Patients with an eligible sinus-rhythm index ECG & 151,941 & 151,941 \\
Meeting the prevalent or concurrent AF criterion & 14,469 & 14,469 \\
No outcome and follow-up below 365 days & 88,945 & 88,945 \\
First outcome after 365 days & 7,481 & not applicable \\
Cases (one-year) or patients with an incident outcome (survival) & 1,855 & 8,775 \\
Controls & 41,770 & 41,770 \\
Total & 43,625 & 50,545 \\
Incident outcomes by 1, 2 and 5 years & 1,855; 2,828; 5,262 & 1,855; 2,828; 5,262 \\
Control follow-up, median (IQR), days & 1,601 (880 to 2,559) & 1,601 (880 to 2,559) \\
Controls with any later ECG & 98.1\% & 98.1\% \\
Controls whose latest later ECG did not read AF & 96.5\% & 96.5\% \\
Cohort file used for modelling & 1,855 cases, 16,000 controls & 4,000 events, 16,000 censored \\
Analysis set & 9,275 (1,855 cases, 7,420 controls; derived) & 20,000 \\
Baseline characteristics by outcome & \textbf{[GAP-T1]} & \textbf{[GAP-T1]} \\
\bottomrule
\end{tabularx}

\smallskip
\raggedright
Note: Exclusion counts are patients meeting each criterion and overlap. IQR, interquartile range. The comparison set size and its 7,420 controls are derived from four controls per case. \textbf{[GAP-T1]} is filled by descriptive statistics of age, sex and the seven ECG measurements by outcome.
\end{table}

\subsection{One-year discrimination}

Table 2 gives the AUROC of the seven models in the one-year comparison set. BTTN + clinical had the highest point AUROC, 0.7776 (95\% CI 0.7664 to 0.7888), and BTTN + SSL + clinical reached 0.7775 (0.7664 to 0.7885). The clinical model reached 0.7597 (0.7488 to 0.7713) and the CNN 0.735 (0.7226 to 0.7469).

Against the clinical model, BTTN + clinical had a higher point AUROC with overlapping 95\% CIs. Both models use the same nine covariates, and BTTN + clinical adds the BTTN log-risk to them, so this comparison holds the covariate set fixed. The two models also differ in loss and fitting, so the difference reflects the ECG pathway together with these differences. Overlap of two marginal CIs on the same patients neither shows nor excludes a paired difference. Against the CNN, the 95\% CIs do not overlap, the lower bound of 0.7664 lying above the CNN's upper bound of 0.7469. This comparison has unequal inputs, since the CNN receives no covariates.

The like-for-like ECG comparison is between ECG-only BTTN and the CNN on the same tensor, folds and patients. ECG-only BTTN reached 0.7397 (0.7275 to 0.7517) against 0.735 for the CNN, with overlapping CIs. BTTN thus reached discrimination comparable to the CNN of this study, with a point AUROC at least as high, while all of its ECG evidence passes through the 12 frequency-by-lead block projections.

Linear BTTR on the same tensor reached 0.5967 (0.5824 to 0.6115), with a CI that does not overlap that of ECG-only BTTN. The trained nonlinear block-term network thus raised discrimination from the value of linear block-term regression to the range of the CNN on the same input. The two block-term models also differ in loss and fitting, so this difference reflects the nonlinear stage together with these differences.

BTTN + SSL reached 0.7414 (0.7292 to 0.7536) and BTTN + SSL + clinical 0.7775, at the level of BTTN and BTTN + clinical. The headline model, BTTN + clinical, uses no pretraining.

ECG-only BTTN had a lower point AUROC than the clinical model, 0.7397 against 0.7597, with CIs overlapping between 0.7488 and 0.7517. Seven of the nine clinical covariates are machine measurements of the same index ECG, so the clinical model is itself a strong ECG-informed comparator.

Paired patient-bootstrap AUROC differences: \textbf{[GAP-C4]}, filled by the mean differences with 95\% CIs for BTTN + clinical minus the clinical model, BTTN + clinical minus the CNN and BTTN minus the CNN, computed from the out-of-fold scores of a run of the one-year comparison on the ICD-coded cohort that writes them to a committed results file.

\begin{table}[htbp]
\centering
\footnotesize
\caption{One-year discrimination of seven models on the ICD-coded comparison set (9,275 patients, 1,855 cases).}
\label{tab:oneyear}
\begin{tabularx}{\textwidth}{@{}L{0.2\textwidth}YYYYYY@{}}
\toprule
Model & ECG input & Clinical covariates & Self-supervised pretraining & AUROC & 95\% CI & Paired difference vs clinical model and vs CNN \\
\midrule
Clinical model (logistic regression) & none & nine & no & 0.7597 & 0.7488 to 0.7713 & \textbf{[GAP-C4]} \\
Linear BTTR & spectral tensor & none & no & 0.5967 & 0.5824 to 0.6115 & \textbf{[GAP-C4]} \\
CNN & spectral tensor & none & no & 0.735 & 0.7226 to 0.7469 & \textbf{[GAP-C4]} \\
BTTN & spectral tensor & none & no & 0.7397 & 0.7275 to 0.7517 & \textbf{[GAP-C4]} \\
BTTN + SSL & spectral tensor & none & yes & 0.7414 & 0.7292 to 0.7536 & \textbf{[GAP-C4]} \\
BTTN + clinical & spectral tensor & nine & no & 0.7776 & 0.7664 to 0.7888 & \textbf{[GAP-C4]} \\
BTTN + SSL + clinical & spectral tensor & nine & yes & 0.7775 & 0.7664 to 0.7885 & \textbf{[GAP-C4]} \\
\bottomrule
\end{tabularx}

\smallskip
\raggedright
Note: AUROC is the mean over 2,000 patient-bootstrap resamples of out-of-fold scores averaged over five seeds of five-fold cross-validation, with percentile 95\% CI. The nine covariates are age, sex, RR interval, P-wave duration, QRS duration, QT interval and the P, QRS and T axes. Model names map to the labels of the released results file as follows: clinical model, ``clinical (age/sex + ECG intervals)''; linear BTTR, ``plain BTTR (spectral)''; CNN, ``CNN (black-box)''; BTTN, ``conv-BTTN (ECG-only)''; BTTN + SSL, ``conv-BTTN + SSL''; BTTN + clinical, ``conv-BTTN + clinical''; BTTN + SSL + clinical, ``conv-BTTN + SSL + clinical (full)''. \textbf{[GAP-C4]} is filled by paired patient-bootstrap AUROC differences from a run of the one-year comparison that writes them to a committed results file.
\end{table}

\subsection{Multi-horizon discrimination from one survival model}

Table 3 reports AUROC at four horizons from one survival BTTN and one BTTN + clinical survival model, each trained once, and from the clinical Cox model. The BTTN + clinical survival model had a higher point AUROC than the clinical Cox model at every horizon: 0.723 (0.701 to 0.745) against 0.71 at 90 days after blanking, 0.727 (0.711 to 0.743) against 0.711 at one year, 0.74 (0.727 to 0.753) against 0.727 at two years, and 0.745 (0.735 to 0.755) against 0.731 at five years. The results file holds no CI for the clinical Cox model. ECG-only survival BTTN ranged from 0.658 to 0.676 across horizons, and adding the clinical term raised the point AUROC at each horizon, to between 0.723 and 0.745.

Tables 2 and 3 use different analysis sets and different clinical comparators. Table 2 is the one-year case-control set of 9,275 patients with a logistic clinical model. Table 3 is the survival set, in which the one-year evaluation includes 19,986 patients with 838 cases, with a penalised clinical Cox model. Their one-year values are therefore not comparable across the two tables, and each table compares models only within its own set.

\begin{table}[htbp]
\centering
\footnotesize
\caption{AUROC at four horizons from one survival model, survival cohort (4,000 events, 16,000 censored).}
\label{tab:horizons}
\begin{tabularx}{\textwidth}{@{}L{0.18\textwidth}YYYYY@{}}
\toprule
Horizon & Patients evaluated & Cases & BTTN (95\% CI) & BTTN + clinical (95\% CI) & Clinical Cox model \\
\midrule
90-day onset after blanking & 20,000 & 410 & 0.658 (0.632 to 0.684) & 0.723 (0.701 to 0.745) & 0.71 \\
1 year & 19,986 & 838 & 0.666 (0.647 to 0.684) & 0.727 (0.711 to 0.743) & 0.711 \\
2 years & 16,982 & 1,299 & 0.674 (0.659 to 0.688) & 0.74 (0.727 to 0.753) & 0.727 \\
5 years & 11,007 & 2,425 & 0.676 (0.665 to 0.687) & 0.745 (0.735 to 0.755) & 0.731 \\
\bottomrule
\end{tabularx}

\smallskip
\raggedright
Note: AUROC is computed on out-of-fold risk from five-fold cross-validation, with percentile 95\% CIs from 2,000 patient-bootstrap resamples for the BTTN models. No CI is available for the clinical Cox model. Cases had an outcome by the horizon, controls were outcome-free beyond it, and patients censored before it were excluded. The 90-day horizon counts outcomes from day 31 to day 120. Values are not comparable with Table 2 (different analysis set and clinical comparator).
\end{table}

\subsection{Calibration, clinical utility and risk stratification}

Figure 2 shows calibration, decision curves and the risk gradient at one year, and Table 4 lists operating points. After cross-fitted Platt scaling, the Brier scores were 0.039 for BTTN + clinical, 0.0396 for BTTN and 0.0398 for the clinical model. With a one-year event rate near 0.04 in this set (Figure 2, right panel), the Brier score is dominated by the event rate, so the scores are read together with the reliability curves and are not used to rank the models.

Flagging the top 10\% of patients by BTTN + clinical risk gave sensitivity 0.272, specificity 0.908, a PPV recalibrated to the cohort incidence of 11.6\% (0.1155) and a number needed to monitor of 8.7. The same flag rate with the clinical model gave sensitivity 0.24, specificity 0.906, PPV 10.2\% (0.1019) and a number needed to monitor of 9.8. The eligible cohort's one-year incidence is 0.0425 (derived from 1,855 / 43,625), so the PPV of BTTN + clinical at this flag rate is about 2.7 times the incidence (derived). Table 4 gives the full grid at 5\%, 10\%, 20\% and 30\%.

Observed one-year AF was 0.55\% in the lowest decile of BTTN + clinical risk and 11.4\% in the highest, a difference of about twentyfold (derived ratio 20.7). In the decision curve, the three models had net benefit above the treat-all and treat-none strategies at low thresholds, and the curves converged towards zero at higher thresholds.

% Image alt text in draft_v5: Figure 2. One-year calibration, decision curve and risk gradient.
\begin{figure}[htbp]
\centering
\includegraphics[width=\textwidth]{../../code/results/clin_eval.png}
\caption{One-year calibration, clinical utility and risk gradient of the survival models after two-fold cross-fitted Platt scaling (source file \texttt{code/results/clin\_eval.png}). Left: Reliability curves by decile of predicted risk for BTTN + clinical, BTTN and the clinical Cox model, with Brier scores in the legend rounded to three decimals. Middle: Decision curves of net benefit across threshold probabilities, with treat-all and treat-none strategies. Right: Observed one-year AF rate by decile of BTTN + clinical risk, with the dashed line at the one-year event rate of the evaluation set.}
\label{fig:clin-eval}
\end{figure}

\begin{table}[htbp]
\centering
\scriptsize
\caption{One-year operating points of the survival models after cross-fitted Platt scaling.}
\label{tab:operating}
\begin{tabularx}{\textwidth}{@{}L{0.16\textwidth}YYYYYYY@{}}
\toprule
Model & AUROC (95\% CI) & Flagged & Sensitivity & Specificity & PPV at cohort incidence & Number needed to monitor & Brier score \\
\midrule
BTTN + clinical & 0.726 (0.710 to 0.742) & 5\% & 0.154 & 0.955 & 0.1307 & 7.7 & 0.039 \\
BTTN + clinical & 0.726 (0.710 to 0.742) & 10\% & 0.272 & 0.908 & 0.1155 & 8.7 & 0.039 \\
BTTN + clinical & 0.726 (0.710 to 0.742) & 20\% & 0.461 & 0.811 & 0.0978 & 10.2 & 0.039 \\
BTTN + clinical & 0.726 (0.710 to 0.742) & 30\% & 0.6 & 0.713 & 0.085 & 11.8 & 0.039 \\
BTTN & 0.666 (0.647 to 0.683) & 5\% & 0.123 & 0.953 & 0.1044 & 9.6 & 0.0396 \\
BTTN & 0.666 (0.647 to 0.683) & 10\% & 0.232 & 0.906 & 0.0983 & 10.2 & 0.0396 \\
BTTN & 0.666 (0.647 to 0.683) & 20\% & 0.387 & 0.808 & 0.0821 & 12.2 & 0.0396 \\
BTTN & 0.666 (0.647 to 0.683) & 30\% & 0.514 & 0.709 & 0.0729 & 13.7 & 0.0396 \\
Clinical Cox model & 0.708 (0.692 to 0.725) & 5\% & 0.129 & 0.953 & 0.1094 & 9.1 & 0.0398 \\
Clinical Cox model & 0.708 (0.692 to 0.725) & 10\% & 0.24 & 0.906 & 0.1019 & 9.8 & 0.0398 \\
Clinical Cox model & 0.708 (0.692 to 0.725) & 20\% & 0.428 & 0.81 & 0.091 & 11.0 & 0.0398 \\
Clinical Cox model & 0.708 (0.692 to 0.725) & 30\% & 0.574 & 0.712 & 0.0813 & 12.3 & 0.0398 \\
\bottomrule
\end{tabularx}

\smallskip
\raggedright
Note: Evaluation set is the one-year at-risk subset of the survival cohort (19,986 patients, 838 cases). AUROC here is the mean over 2,000 patient-bootstrap resamples of the calibrated probabilities and differs slightly from Table 3. Flagged is the share of the evaluation set flagged as high risk. PPV and number needed to monitor are recomputed from sensitivity and specificity at the eligible cohort's one-year incidence, 1,855 / (1,855 + 41,770). Decile rates in the text are rounded to 0.55\% and 11.4\%.
\end{table}

\subsection{Faithfulness of the block signatures}

The block scores that the probes read are computed from the block-term layer alone and exclude the temporal stage, the head and the clinical term by construction (Sections 3.4 and 3.5). The five analyses of Section 3.8 are held at the placeholders below until their ICD-coded results are available.

\begin{itemize}
\item Control-probe results: \textbf{[GAP-F1]} (Table 5), filled by the AUROCs, CIs and recoveries of the five probe sources from a run whose log records the 9,275-patient ICD-coded comparison set.
\item Intervention curve: \textbf{[GAP-F2]} (Figure 4), filled by the top-n and bottom-n removal curves from a run on the ICD-coded comparison set.
\item Block-structure ablation: \textbf{[GAP-F3]} (Table 6), filled by full AUROC and probe recovery of the block-term, CP and full-rank variants from a run on the ICD-coded comparison set.
\item Block-count sweep: \textbf{[GAP-F4]} (Table 6), filled by the mean and standard deviation of AUROC for K = 4, 8, 12, 16 and 24 from a run on the ICD-coded comparison set.
\item Substrate map: \textbf{[GAP-F5]} (Figure 3), filled by the frequency and lead occlusion importances and the mean block signature from a run on the ICD-coded comparison set.
\end{itemize}

\refstepcounter{table}
\gapbox{\textbf{[GAP-F1]} \textbf{Table 5.} Control probes: AUROC (95\% CI) and recovery for probes on the block scores, the raw flattened tensor, the CNN penultimate features, a matched-dimension random projection and the label-blind block scores. Filled by a run of the control-probe analysis on the ICD-coded comparison set.}
\label{tab:probes}

\refstepcounter{table}
\gapbox{\textbf{[GAP-F3, GAP-F4]} \textbf{Table 6.} Block-structure ablation (block-term, CP and full-rank first layers: full AUROC with 95\% CI and probe recovery) and block-count sweep (K = 4, 8, 12, 16, 24: mean and standard deviation of AUROC). Filled by the block-structure ablation and block-count sweep outputs of runs on the ICD-coded comparison set.}
\label{tab:ablation}

\refstepcounter{figure}
\gapbox{\textbf{[GAP-F5]} \textbf{Figure 3.} Substrate map: occlusion importance per frequency bin and per lead and the mean absolute block signature. Filled by the substrate-map output of a run on the ICD-coded comparison set.}
\label{fig:substrate}

\refstepcounter{figure}
\gapbox{\textbf{[GAP-F2]} \textbf{Figure 4.} Intervention curve: full-model AUROC after cumulative removal of the top-n and bottom-n ranked frequency bins and leads. Filled by a plot of the intervention-curve output of a run on the ICD-coded comparison set.}
\label{fig:intervention}

\subsection{External analyses}

The external analyses of Section 3.9 address whether the method and its frequency and lead signatures transfer to CODE-15~\cite{ribeiro2021code15} and PTB-XL~\cite{wagner2020}. Neither dataset has dated outcomes, so these analyses do not test incident prediction, and the results of Sections 4.2 to 4.4 do not depend on them.

\begin{itemize}
\item External analyses: \textbf{[GAP-X1]} (Table 7 and a supplementary figure), filled by the five-fold cross-validated AUROCs for CODE-15 identification, CODE-15 detection and PTB-XL detection and by their frequency and lead occlusion importances, from runs of the external analyses that write these values to a committed results file.
\end{itemize}

The external analyses complement the internal validation by showing whether the method transfers to independently collected ECGs. Table 7 lists the task, dataset, patients, cases and cross-validated AUROC of each analysis (\textbf{[GAP-X1]}).

\refstepcounter{table}
\gapbox{\textbf{[GAP-X1]} \textbf{Table 7.} External analyses: task, dataset, patients and cases, and five-fold cross-validated AUROC for CODE-15 identification, CODE-15 detection and PTB-XL detection. Filled by runs of the external analyses that write the cross-validated AUROCs to a committed results file.}
\label{tab:external}

\section{Discussion}

\subsection{Principal findings}

BTTN + clinical reached the highest point AUROC of the seven models compared for one-year incident AF on ICD-coded MIMIC-IV-ECG outcomes, 0.7776 (0.7664 to 0.7888), against 0.7597 for the clinical model and 0.735 for the CNN. Its CI overlaps that of the clinical model and lies above that of the CNN, and the paired differences are held at \textbf{[GAP-C4]}, filled by a one-year comparison run that writes them to a results file. ECG-only BTTN and the CNN of this study reached comparable discrimination from the same tensor, 0.7397 and 0.735, so BTTN reached a point AUROC at least as high as this CNN while its ECG evidence passes through readable frequency and lead projections. One BTTN + clinical survival model gave a higher point AUROC than the clinical Cox model at each of four horizons, from 0.723 at 90 days to 0.745 at five years. After recalibration its one-year Brier score was 0.039, and observed one-year AF was 0.55\% in the lowest decile of its risk and 11.4\% in the highest. The faithfulness of the signatures is assessed by five analyses (\textbf{[GAP-F1 to GAP-F5]}, filled by their results on the ICD-coded comparison set).

\subsection{Comparison with AI-ECG and clinical risk models}

The results follow the complementarity of ECG and clinical information reported earlier. Khurshid and colleagues found the combination of CHARGE-AF and an ECG model above either alone in three cohorts~\cite{khurshid2022}, Christopoulos and colleagues reported 0.72 for the combination against 0.69 for each component~\cite{christopoulos2020}, and Biton and colleagues obtained the highest AUROC by combining ECG features, learned representations and clinical data~\cite{biton2021}. Here the same direction appears at one year and at every survival horizon. The clinical model of this study, a strong comparator as in prior work~\cite{alonso2013}, includes seven measurements of the index ECG, among them P-wave duration~\cite{nielsen2015}, and the contribution of BTTN shows in fusion with the same covariates. Interpretable hand-crafted ECG features have predicted AF minutes ahead in a small database~\cite{hammer2022}, whereas BTTN carries readable ECG evidence to horizons of one to five years.

Published AUROCs differ in task and horizon: new-onset AF within one year (0.85)~\cite{raghunath2021}, five-year AF or flutter with MIMIC-IV as an external set (0.77)~\cite{jabbour2024}, AF within 31 days (0.86)~\cite{yuan2023}, and AF identified during sinus rhythm (0.87)~\cite{attia2019}. Cohorts, label sources and exclusion windows differ, so no ranking across studies is implied. The outcome here is a dated hospital diagnosis code, whereas foundation models have relied on labels extracted from ECG reports~\cite{nolinlapalme2026} and the abstract of another does not state its AF label source in MIMIC-IV~\cite{ko2026}. On MIMIC-IV-ECG, AF detection has reached an AUROC of 0.996~\cite{mckeen2025}, close to the ceiling, and incident prediction after a 30-day blanking window is a harder task.

BTTN + SSL and BTTN + SSL + clinical reached 0.7414 and 0.7775, at the level of their supervised counterparts. Self-supervised gains on 12-lead ECG reported elsewhere were modest: the best self-supervised method came within 0.5\% of supervised training in linear evaluation and fine-tuning added about 1\%~\cite{mehari2022}, and a self-supervised foundation model reached 0.742 for five-year AF risk against 0.734 for its supervised counterpart~\cite{nolinlapalme2026}.

\subsection{Interpretability by design and its measurement}

In BTTN, every piece of ECG evidence reaches the temporal stage through the projections onto $A_k$ and $C_k$, and these factors are model parameters that can be read directly as frequency and lead profiles, through the six rank-one maps of each block or the two-channel signature (Section 3.4). This places BTTN with interpretable-by-design models~\cite{rudin2019,koh2020,bohle2022}. It also sets BTTN apart from explanations computed after training, such as LIME, Grad-CAM and SHAP in ECG studies~\cite{hughes2021,storas2023,matsuo2026}. The clinical term adds age, sex and ECG measurements as a readable weight vector without entering the block layer, so fusion keeps the structure of the ECG pathway.

The claim made by construction is limited to which frequency and lead directions enter the model. The temporal stage and head are nonlinear and could recombine block channels in ways the factors alone do not show, which is the leakage concern raised for concept models~\cite{mahinpei2021,havasi2022}. Ghassemi and colleagues argued that explanations of this kind risk false reassurance and that validation should take their place~\cite{ghassemi2021}. This paper answers by measurement. The control probes, the intervention test and the ablations of Section 3.8 quantify how far the readouts track the model, and their interpretation is held at \textbf{[GAP-F1 to GAP-F5]}, filled by those results on the ICD-coded comparison set.

The substrate map is read against established sinus-rhythm markers of AF risk, P-wave duration~\cite{nielsen2015} and P-wave terminal force in lead V1~\cite{huang2020,wolder2023}. That comparison is \textbf{[GAP-F5]}, filled by the substrate map from the ICD-coded comparison set read against these markers. Frequency analysis of the surface ECG is established for atrial activity during AF, where the atrial fibrillatory rate relates to atrial cycle length and AF persistence~\cite{bollmann2006,platonov2014}, and whether any block signature in sinus rhythm relates to these measures is also \textbf{[GAP-F5]}.

\subsection{Clinical implications}

Targeted monitoring needs a score that concentrates future cases among few patients. Flagging the top 10\% by BTTN + clinical risk captured 27.2\% of one-year cases with a PPV of 11.6\% at the cohort incidence, so about nine patients would be monitored per future case (number needed to monitor 8.7, against 9.8 with the clinical model). Noseworthy and colleagues found that AI-ECG risk stratification raised AF detection in high-risk patients during up to 30 days of monitoring~\cite{noseworthy2022}, and Raghunath and colleagues reported a number needed to screen of 9 in a simulated deployment~\cite{raghunath2021}. Those figures come from other cohorts and definitions and serve as context only.

The difference between the lowest risk decile (0.55\%) and the highest (11.4\%) indicates that monitoring intensity could be graded by risk. A single survival model gives horizon-specific risk at 90 days and at one, two and five years (Table 3), so a monitoring programme can choose its horizon without retraining. Because the clinical contribution is a weight vector, the covariates that raise a given patient's clinical term can be read off, and the ECG contribution can be traced to block factors whose frequency and lead content is visible.

\subsection{Limitations}

All data come from one health system~\cite{johnson2023}, and no external cohort with dated outcomes was available. The external analyses address identification and detection, a different task, with in-sample importances, and their results are \textbf{[GAP-X1]}, filled by the external analysis runs. External dated-outcome validation, then prospective evaluation, is the next study needed.

Outcomes are hospital codes dated by admission. AF diagnosed only in outpatient care or never diagnosed is missed, a coded date can lag the true onset, and code I48 includes atrial flutter, as does the outcome of Jabbour and colleagues~\cite{jabbour2024}. The blanking window and the exclusion of patients with an AF-reading ECG reduce misclassification of prevalent AF as incident without removing it.

Control follow-up is inferred from the last ECG or death. Of the 41,770 one-year controls, 40,960 (98.1\%) had a later ECG and 40,306 (96.5\%) had a latest later ECG that did not read AF, so 654 controls (1.6\%, derived) had a latest ECG reading AF without a hospital AF code, which the coded outcome misses.

Two covariate approximations (Section 3.2) affect the clinical model, the clinical term and the clinical Cox model alike. Age is the anchor age rather than age at the index ECG, and the unrecoded missing codes leave 819 rows of the one-year cohort file with a P-wave duration near 29,600 ms. Removing both requires rerunning the clinical, fused and multi-horizon analyses.

The CNN comparator is a compact network on the spectral tensor, smaller than the architectures trained on raw waveforms in published AI-ECG studies, so the ECG-only comparison is with this CNN.

Case-control sampling with four controls per case leaves AUROC, a rank statistic~\cite{hanley1982}, comparable across the seven models, while PPV and number needed to monitor were recomputed at the cohort incidence and not measured.

Comparisons rest on point estimates and marginal CIs until the paired differences from a one-year comparison run fill \textbf{[GAP-C4]}. Discrimination over time is reported as horizon AUROCs, without a concordance index.

Until \textbf{[GAP-F1 to GAP-F5]} are filled by runs of the five analyses on the ICD-coded cohort, the interpretability claim covers the structure of the model only.

The 30,000-ECG pretraining pool may include cohort patients, which affects only the SSL rows, whose pretraining used no labels.

The 50 Hz notch was applied to recordings from a United States hospital with 60 Hz mains. Both frequencies lie above the 40 Hz tensor limit, and spectral leakage into the highest retained bins was not measured.

Baseline characteristics are \textbf{[GAP-T1]}, filled by descriptive statistics of the cohort files.

\section{Conclusion}

A block-term tensor network that routes all ECG evidence through 12 readable frequency-by-lead block projections, and adds age, sex and seven ECG measurements as a separate linear clinical term, reached the highest point AUROC of seven models compared for one-year incident AF from one sinus-rhythm 12-lead ECG on ICD-coded MIMIC-IV outcomes. Its AUROC was 0.7776 (95\% CI 0.7664 to 0.7888), against 0.7597 (0.7488 to 0.7713) for a clinical risk-factor model and 0.735 (0.7226 to 0.7469) for the CNN of this study on the same tensor, folds and patients. Without the clinical term, BTTN reached discrimination comparable to this CNN. Trained once as a survival model, BTTN + clinical gave a higher point AUROC than a clinical Cox model at 90 days and at one, two and five years, and observed one-year AF was 0.55\% in the lowest decile of its risk and 11.4\% in the highest. Against the CNN these point estimates have non-overlapping CIs. Against the clinical model the CIs overlap, and the comparison is completed by the paired differences at \textbf{[GAP-C4]}, filled by a run of the one-year comparison that writes them to a results file.

The interpretability of BTTN rests on its structure. The frequency and lead directions that the model can use are its block factors, the clinical contribution is a weight vector outside the block layer, and the block scores can be read without the nonlinear head. The faithfulness of the signatures is assessed by five analyses (\textbf{[GAP-F1 to GAP-F5]}), filled when the control probes, the intervention curve, the two ablations and the substrate map are available on the ICD-coded cohort. Together with validation in an external cohort with dated outcomes, those analyses decide how far the readable signatures can support clinical interpretation of an individual risk estimate.

\section*{Declarations}

\textbf{Data availability.} MIMIC-IV-ECG version 1.0~\cite{gow2023} (doi:10.13026/4nqg-sb35) and the MIMIC-IV hospital module~\cite{johnson2023} are available under credentialed access on PhysioNet~\cite{goldberger2000}. CODE-15~\cite{ribeiro2021code15} and PTB-XL~\cite{wagner2020} were used for the external analyses. No raw data are redistributed. The analysis code is released under the MIT license. Repository URL or archive DOI: \textbf{[GAP-REPO]}, filled by the public repository URL or archive DOI of the released code.

\textbf{Ethics statement.} \textbf{[GAP-ETH]}, filled by the authors' wording on ethics approval or exemption, waiver of consent, de-identification of the data and the data use agreements for MIMIC-IV-ECG and the MIMIC-IV hospital module.

\textbf{Author contributions (CRediT).} \textbf{[GAP-CREDIT]}, filled by the CRediT roles of each author.

\textbf{Conflicts of interest.} \textbf{[GAP-COI]}, filled by each author's declaration of competing interests.

\textbf{Funding.} \textbf{[GAP-FUND]}, filled by the funders and grant numbers, or a statement that no specific funding was received.

\textbf{Declaration of generative AI use.} During the preparation of this work the authors used Claude (Anthropic), model Claude Opus 5.5, through the agents of the academic-paper skill, for literature search and verification of the cited references, outlining, argument planning, drafting of the manuscript, internal review and revision of the draft, citation checking, drafting of the abstract, keywords and highlights, and formatting. The agents also counted rows and computed derived counts and ratios from the authors' committed result files, each marked as derived in the text; they trained no model and produced none of the reported model results. No AI tool is listed as an author. Authors' verification statement: \textbf{[GAP-AIUSE]}, filled by the named authors' statement that they reviewed and edited the content and take responsibility for it.

\textbf{Supplementary material.} TRIPOD+AI checklist~\cite{collins2024}: \textbf{[GAP-TRIPOD]}, filled by the TRIPOD+AI checklist that the authors complete against the final manuscript. Table S1, settings of preprocessing, BTTN, the clinical term, SSL pretraining and the comparators, given below. Supplementary figure of external frequency and lead importances: \textbf{[GAP-X1]}, filled by runs of the external analyses that write the importances to a results file.

\renewcommand{\thetable}{S1}
\begin{table}[htbp]
\centering
\small
\caption{Settings of preprocessing, models and evaluation, as fixed in the released analysis code.}
\label{tab:S1}
\begin{tabularx}{\textwidth}{@{}L{0.25\textwidth}Y@{}}
\toprule
Component & Setting \\
\midrule
High-pass filter & Butterworth, order 2, 0.5 Hz, forward and backward \\
Notch filter & 50 Hz, quality factor 30, forward and backward \\
Wavelet denoising & sym4, level 3, soft threshold of 7 times the standard deviation of the finest detail coefficients \\
Resampling (MIMIC-IV-ECG) & decimation by 2, 500 to 250 Hz, zero phase \\
Spectrogram & magnitude, 128-sample window, 32-sample hop, frequencies up to 40 Hz, log(1 + x), 40 time frames, per-recording standardisation \\
BTTN block-term layer & K = 12 blocks, frequency factor rank 3, lead factor rank 2, 72 channels, no core tensor \\
BTTN temporal stage & grouped convolutions, kernel sizes 3, 7 and 15, 6 filters per channel and kernel size, batch normalisation, rectified linear unit, softmax attention and max pooling \\
BTTN head & batch normalisation, dropout 0.4, multilayer perceptron with 256 and 128 hidden units \\
BTTN training & within-batch Cox partial likelihood, Adam, learning rate 2 x 10$^{-3}$, weight decay 10$^{-4}$, cosine schedule, batch 512, at most 120 epochs, 85/15 inner split, early stopping on validation AUROC with patience 16, three-model ensemble \\
Clinical term & linear, no bias, nine standardised covariates, warm start from a class-balanced logistic regression, no weight decay \\
SSL pretraining & NT-Xent, temperature 0.2, 30,000 ECGs one per patient, 60 epochs, Adam, learning rate 10$^{-3}$, weight decay 10$^{-5}$, cosine schedule, batch 256, projection head of 256 and 128 units \\
SSL views & masks of 20\% of time frames and 20\% of frequency bins, lead dropout probability 0.2, Gaussian noise of standard deviation 0.05, roll of up to 3 frames \\
CNN & convolutions of 32, 64 and 64 filters, kernel sizes 7, 5 and 3, batch normalisation, rectified linear unit, max pooling after the first two, average pooling after the third, dropout 0.3, 0.4 and 0.5, per-channel standardisation \\
CNN training & binary cross-entropy with positive-class weight equal to the control-to-case ratio, Adam, learning rate 10$^{-3}$, weight decay 10$^{-3}$, batch 256, at most 80 epochs, early stopping on validation AUROC with patience 12, one 85/15 inner split, three-model ensemble \\
Linear BTTR & up to 8 blocks, rank-one enforced, signal-to-noise grid 20 to 40 dB in steps of 5, ratios 95, 97 and 99, labels coded -1 and +1 \\
Clinical model & median imputation, standardisation, logistic regression with the library default L2 penalty, balanced class weights \\
Clinical Cox model & Cox proportional hazards with L2 penalty 0.1, median imputation \\
One-year comparison & four controls per case, stratified five-fold cross-validation repeated with five seeds, 2,000 patient-bootstrap resamples \\
Multi-horizon analysis & one five-fold run stratified by event, early stopping on validation AUROC for an event within one year \\
\bottomrule
\end{tabularx}
\end{table}

\bibliographystyle{vancouver}
\bibliography{references}

\end{document}
